\begin{table}[htbp]
\centering
\caption{Border $\times$ Restricted interaction: New Mexico vs. Colorado}
\label{tab:regime_interaction}
\footnotesize
\begin{tabular}{lc}
\toprule
Outcome & NM+CO \\
\midrule
Revenue per capita & $-0.458$ \\
 & (0.388) \\
Dispensary density & $-0.932^{*}$ \\
 & (0.488) \\
Revenue per store & $+0.465$ \\
 & (0.314) \\
\bottomrule
\end{tabular}
\begin{minipage}{0.98\textwidth}
\vspace{0.5em}

\fontsize{9pt}{11pt}\selectfont
\textit{Notes:} $\log(y_{cst}) = \alpha + \beta\,\text{Border}_c
+ \theta\,(\text{Border}_c \times \text{Restricted}_s)
+ \delta_{st} + \varepsilon_{cst}$.
Restricted$_s$ = 1 for Colorado and 0 for New Mexico. Colorado is restricted to
2022Q2 onward, matching the window we observe for New Mexico, so that every
Colorado county-quarter in the sample is one in which New Mexico was already
legal and Colorado's border classification is unambiguous
(Section~\ref{sec:crossstate}). Both Border$_c$ and Restricted$_s$ are
time-invariant in this sample, so county fixed effects cannot be added because
they would absorb the interaction. State $\times$ quarter fixed effects are
included using a concatenated state-quarter key, and standard errors are
clustered at the county level. Because the entry regime varies only across two
states, county clustering addresses within-county serial correlation but cannot
rule out other persistent state-level differences correlated with the regulatory
regime. We therefore interpret $\theta$ as the difference in the estimated
border premium between Colorado and New Mexico, rather than as an unqualified
causal treatment effect. $^{*}p<0.10$, $^{**}p<0.05$, $^{***}p<0.01$.
\end{minipage}
\end{table}
